\chapter{交流电}\label{ch:g9:alternating-current}

你做过的每一个电路都由电池供电：电流
平稳地从 $+$ 流向 $-$，如钟表的指针般尽职。
但墙上插座输出的电流却更为奇特——一种
每秒换向一百次的电流，
其电压不停地在正负之间摆动。认识一下
交流电：世界上几乎全部电能
都以这种形式产生、输送并供应给用户。

\section{两类电流}

\begin{definition}[直流电与交流电]\label{def:g9:alternating-current:dcac}
称\emph{直流电}\index{直流电}（DC）为
始终沿同一方向流动、由稳定电压驱动的电流，
即电池提供的电流，电压恒定，直到电池耗尽。
\emph{交流电}\index{交流电}（AC）的方向
周期性地反转，由在正负之间平滑
往复变化的电压驱动，周而复始；
插座中的电流即属此类，交流发电机产生的
也是它（下一章将说明其来源）。
\end{definition}

\begin{omfigure}
\begin{tikzpicture}
\begin{axis}[omaxis, width=10.8cm, height=5.2cm,
    xlabel={时间（\unit{ms}）}, ylabel={电压（\unit{V}）},
    xtick={0,10,20,30,40}, ytick={-325,0,230,325},
    yticklabels={$-325$,$0$,$230$,$325$},
    ymin=-400, ymax=400, xmin=0, xmax=42, samples=200]
  \addplot[very thick, omDef, domain=0:40] {230*0+4.5};
  \addplot[very thick, omThm, domain=0:40]
    {325*sin(deg(2*3.14159*0.05*x))};
  \node[right, font=\footnotesize, omDef] at (axis cs:29,42)
    {电池：稳定的 \qty{4.5}{V}};
  \node[right, font=\footnotesize, omThm] at (axis cs:1.2,352)
    {市电：来回摆动};
\end{axis}
\end{tikzpicture}

{\small 两种电压随时间的变化：电池的平直线，
稳定不变；以及市电在正负之间
无尽地摆动，每个完整周期为二十毫秒。}
\end{omfigure}

\begin{definition}[周期与频率]\label{def:g9:alternating-current:period}
交变电压的\emph{周期}\index{周期} $T$
是一个完整周期的时长——电压先上升、
再下降、再回到原处所用的时间，单位为秒。其\emph{频率} $f$ 是
每秒内完成的周期数，单位为赫兹——这也是声音
的单位，此处再次出现——二者互为倒数：
\[
f = \frac{1}{T}, \qquad T = \frac{1}{f}.
\]
\end{definition}

\begin{proposition}[市电的基本参数]\label{prop:g9:alternating-current:mains}
世界上大多数地区的市电以
$f = \qty{50}{Hz}$ 的频率交变——每秒五十个完整周期，周期
为 $T = 1/50 = \qty{0.02}{s}$，即每次摆动二十毫秒
（部分地区为六十赫兹）。台灯电线中的电流
每秒换向一百次，
——每周期换向两次——从未平稳地
从一个接线柱走向另一个。
\end{proposition}

\begin{example}[频率计算]\label{ex:g9:alternating-current:arithmetic}
倒数关系两种算法都成立。一台以
\qty{60}{Hz} 循环运行的发电机：$T = 1/60 \approx \qty{0.017}{s}$——
十七毫秒。示波器显示一个周期持续
\qty{0.04}{s}：$f = 1/0.04 = \qty{25}{Hz}$。一台自行车发电机
在缓慢蹬踩时产生 \qty{10}{Hz}：周期为十分之一
秒——灯泡会明显闪烁，每秒十次脉冲，
直到速度快到肉眼察觉不出闪烁为止。
\end{example}

\section{读取波形}

\begin{method}[读取示波器波形]\label{met:g9:alternating-current:scope}
\emph{示波器}\index{示波器}——或任何
分析软件——在方格屏幕上画出电压随时间的变化；
每格所代表的伏特数与毫秒数
由两个旋钮设定。读取波形的步骤如下：
\begin{enumerate}
  \item 数出一个完整周期所占的格数，再乘以
        每格代表的时间，即 $T$；然后 $f = 1/T$；
  \item 数出从中线到波峰所占的格数，
        再乘以每格代表的伏特数，得到\emph{峰值
        电压} $U_{\text{max}}$；
  \item 检查波形形状：市电与交流发电机
        产生的，正是图中那种平滑且不断重复的波形——
        即\emph{正弦曲线}，是数学中对
        这种完美摆动的称呼。
\end{enumerate}
\end{method}

\begin{proposition}[有效电压]\label{prop:g9:alternating-current:effective}
对于每秒经过从 $+U_{\text{max}}$ 到 $-U_{\text{max}}$ 之间
各数值一百次的摆动，其“电压”究竟是多少？
较为确切的单一数值是\emph{有效
电压}\index{有效电压} $U$：即能使同一电阻器产生相同热效应的稳定直流电压。
对正弦曲线而言，它
远低于峰值——$U_{\text{max}}$ 约为 $1.4$
倍 $U$，而每一只交流电压表设计时所
要读取的正是这个值。插座上著名的“\qty{230}{V}”
是有效值；峰值实际约为
\qty{325}{V}，每秒出现五十次，正、负两个方向各一次。
\end{proposition}

\begin{example}[哪些器件在意电流方向]\label{ex:g9:alternating-current:whocares}
热效应对两种电流都适用——灯丝或
电热水壶的加热丝无论电流朝哪个方向流动都会发热，这正是
灯具和取暖器能毫无怨言地使用交流电的原因。但是
电子设备——手机、笔记本电脑以及一切带芯片的器件——
需要稳定的直流电：它们的充电器和内置的电源
会把插座中摆动的交流电转换为类似电池那样
稳定的直流（摸一摸充电器的温度：转换本身要
付出代价）。电池同样无法用摆动的电流充电：
充电相当于反向进行的化学反应，而化学反应
要求电流方向保持一致。
\end{example}

\begin{remark}[为什么世界选择了摆动的电流]\label{rem:g9:alternating-current:why}
既然电池如此令人安心地稳定，为何文明社会
采用交流电供电？原因有两个，都留待后两章再讲：
交流电恰恰是旋转机械\emph{自然
产生} --- 而且，决定性的是，只有交流电能供给变压器，
这个不起眼、嗡嗡作响的小箱子，使电能能远距离输送至全国各地，
以极高电压输送，进入家中时降为温和的电压。
电网采用交流电，是输送的需要；接下来的
两章将追踪电流的路径，从旋转的线圈
一直到厨房墙上的插座。
\end{remark}

\begin{remark}[摆动的电流与人体]\label{rem:g9:alternating-current:safety}
安全问题的旧算术又多了一个严峻的细节：市电的
每秒 50 次的摆动，接近最容易扰乱心脏
自身电节律的频率 --- 交变电击
在安培数相同时，比稳定电流更危险。
童年时的每条规则依然适用，只是加倍了：插座的
\qty{230}{V}为有效值 --- 峰值为\qty{325}{V} ---
并不仅仅是一块更大的电池，
它是另一种更不饶人的东西。
\end{remark}

\section{练习}

\begin{exercise}[$\star$]\label{exo:g9:alternating-current:1}
各用一句话区分直流电与交流电，并分别说明
各自在家庭中的来源。
\end{exercise}

\begin{exercise}[$\star$]\label{exo:g9:alternating-current:2}
给出周期与频率的定义，并写出二者的
倒数关系。求 \qty{50}{Hz} 时的 $T$，以及 $f$ 在
$T = \qty{0.01}{s}$ 时的值。
\end{exercise}

\begin{exercise}[$\star$]\label{exo:g9:alternating-current:3}
市电电流每秒改变
方向多少次，取 \qty{50}{Hz}？（注意：每个周期内的换向次数。）
\end{exercise}

\begin{exercise}[$\star$]\label{exo:g9:alternating-current:4}
示波器上一个周期横跨 $4$ 格，
每格 \qty{5}{ms}。求 $T$ 与 $f$。
\end{exercise}

\begin{exercise}[$\star$]\label{exo:g9:alternating-current:5}
插座标示 \qty{230}{V}；峰值达到
\qty{325}{V}。说出这两个数值，并解释
较小那一个的含义。
\end{exercise}

\begin{exercise}[$\star$]\label{exo:g9:alternating-current:6}
为什么电热水壶欣然接受交流电，而手机充电器必须
将其转换？试用“一种效应、一种化学”来说明。
\end{exercise}

\begin{exercise}[$\star$]\label{exo:g9:alternating-current:7}
画出（或描述）两种波形：一节 \qty{9}{V} 电池，
以及周期为 \qty{20}{ms}、峰值为
\qty{9}{V}的交变电压。仅看示波器所画的图形，有哪一个单一特征
能一眼将它们区分开？
\end{exercise}

\begin{exercise}[$\star$]\label{exo:g9:alternating-current:8}
给出文明世界采用交流电供电的两个理由，
即本章预告的两点。
\end{exercise}

\begin{exercise}[$\star\star$]\label{exo:g9:alternating-current:9}
速度慢时，自行车发电机点亮的灯会明显闪烁；
蹬得越快，灯光越稳定。请用频率解释这两种情况，
并估算频率低于多少时闪烁
会令人烦扰，已知电影约以
每秒 $24$ 幅画面骗过人眼。
\end{exercise}

\begin{exercise}[$\star\star$]\label{exo:g9:alternating-current:10}
并联在插座两端的交流电压表读数为 \qty{230}{V}。
（设想中的）无畏示波器并联在同一插座两端，
显示的峰值是多少 --- 每秒又有多少个
峰值（正负两个方向都计入）？
\end{exercise}

\begin{exercise}[$\star\star$]\label{exo:g9:alternating-current:11}
某地区电网以 \qty{60}{Hz} 运行，有效值为
\qty{120}{V}。求其周期与近似峰值电压，
并指出你的哪一件旅行用电器会
提出异议（见
\cref{ex:g9:alternating-current:whocares}；现代充电器会细读标签上的小字
）。
\end{exercise}

\begin{exercise}[$\star\star\star$]\label{exo:g9:alternating-current:12}
老式影片中，屏幕与荧光灯有时会出现
滚动的暗条纹 --- 相机的画面
恰好捕捉到灯光亮度摆动途中的瞬间。以 \qty{50}{Hz}
市电供电的灯在每个峰值处都会变亮，不论\emph{哪一种符号}：每秒有多少
次亮度脉冲？相机每秒 $25$ 帧
对这种脉动进行采样：请说明两种节律的
拍频为何会形成缓慢移动的条纹，以及在
\qty{60}{Hz} 的地区拍摄时，图案为何会改变。
\end{exercise}

\section{问题：示波器实验台}

\begin{problem}[{周末问题 --- 示波器实验台上的资格之夜；
四个神秘电源，一块屏幕；
技术员的记录簿}]\label{pb:g9:alternating-current:1}
须根据波形辨认四个未贴标签的电源。
示波器旋钮设定为：水平方向每格 \qty{5}{ms}，
竖直方向每格 \qty{2}{V}（除另有说明处外
）。由你负责保管记录簿。

\medskip\noindent\textbf{第 I 部分 --- 四条波形。}
\begin{enumerate}
  \item 电源 A 为一条平直线，高于中线 $2.25$ 格；
        判断其种类与电压，
        并指出它最可能对应哪种家用电源。
  \item 电源 B 为平滑波形，每 $4$ 格一个周期，
        波峰高 $3$ 格。说明其种类、周期、
        频率与峰值电压？
  \item 电源 C 为一条平直线，\emph{低于}中线
        $2.25$ 格。它是什么？是哪一个微小误差
        造成了这一现象？
  \item 电源 D 为平滑波形，每 $2$ 格一个周期，
        波峰高 $1.5$ 格。求其全部特征量，
        并说明：若两者接同一盏灯，
        D 的灯为何比 B 的灯闪烁得更少？
\end{enumerate}

\medskip\noindent\textbf{第 II 部分 --- 技术员的
问题。}
\begin{enumerate}[resume]
  \item 并联在电源 B 两端的交流电压表示数约为
        \qty{4.2}{V}，而其波峰达到 \qty{6}{V}。
        请解释这两个数值为何并不矛盾。
  \item 四个电源中，哪一个可以直接为手机充电，
        而其余电源与手机电池之间须接入
        何种元件？
  \item 台面上的警示牌写道：“同样的伏特数下，
        交流电更伤人。”请复述其原因。
  \item 一名学员问：为何示波器所画的 B 波形
        如此平滑圆润而非锯齿状？请说明该波形
        的名称及其来源（哪一种机器，
        在下一章讲到，会自然地产生它？）
\end{enumerate}

\medskip\noindent\textbf{第 III 部分 --- 电网自身的
波形。}
一台经安全隔离的演示装置显示市电波形，
刻度为每格 \qty{100}{V} 与每格 \qty{5}{ms}。
\begin{enumerate}[resume]
  \item 预测该波形：每周期占几格，
        以及波峰高几格。
  \item 由该波形求出 $f$ 与有效电压，
        并与插座上的著名标注对照，
        检验二者是否相符。
  \item 将演示装置改接为电池组，其有效发热功率
        与原来相同。请用一句话描述新的
        波形。
  \item 以技术员的三句话结束记录簿：
        平直线意味着什么，波形意味着什么，
        以及哪一个单一数值使二者可以
        公平比较。
\end{enumerate}
\end{problem}
